\clearpage
\section{T3 in full: all twenty-six rows}

This is the whole table, generated from the stored results rather than transcribed. The knowledge
condition is the recovering set, the estimates are at $n = 20\,000$ with the first seed, and the
population interval behind each band is in the source file. Three declared queries are excluded and
they are listed under the table with the reason.

\begin{landscape}
\begin{table}[p]
\centering
\footnotesize
\setlength{\tabcolsep}{3pt}
\caption*{\textbf{T3.} Sensitivity reports on twenty-six named analyses. $|K|$ is the number of claims the analyst asserts, with the number inside the treatment's own component in parentheses. The two interval columns are the union of the estimate over every knowledge state within one and two elementary revisions of those claims. $r_0$ is the null radius, the number of revisions after which the interval contains zero, from the population, from the band at the first seed, and as the range over twenty seeds. $r_{\mathrm{val}}$ is the breakdown radius, the number of revisions after which the committed adjustment set stops being valid. The last column names the claims whose retraction first opens the interval to zero, in the source's own variable names. Knowledge: the recovering set; estimates at $n = 20\,000$, first seed.}
\begin{tabular}{@{}llccl cc l c l@{}}
\toprule
network & query & $|K|$ & true & estimate [95\,\%] & one revision & two revisions
 & $r_0$ pop / seed / range & $r_{\mathrm{val}}$ & first breaking claims \\
\midrule
\multicolumn{10}{@{}l}{\itshape Published analyses: the treatment and the outcome are the ones the source paper declares.}\\[1pt]
Acid\_1996 & x3 $\to$ x15 & 1 (0) & $-2.169$ & $-2.176$ [$-2.203$, $-2.149$] & [$-2.226$, $-2.147$] & [$-2.226$, $-2.147$] & never / $>$3 / $>$3 & never & -- \\
Didelez\_2010 & HRT $\to$ TCI & 1 (0) & $-0.586$ & $-0.586$ [$-0.599$, $-0.572$] & [$-0.604$, $-0.568$] & [$-0.604$, $-0.568$] & never / $>$3 / $>$3 & never & -- \\
Kampen\_2014 & SUS $\to$ EGC & 3 (0) & $-0.643$ & $-0.647$ [$-0.663$, $-0.630$] & [$-0.663$, $-0.630$] & [$-0.663$, $-0.630$] & never / $>$3 / $>$3 & never & -- \\
Polzer\_2012 & ToothLoss $\to$ Mortality & 5 (0) & $\phantom{-}1.193$ & $\phantom{-}1.190$ [$\phantom{-}1.183$, $\phantom{-}1.197$] & [$\phantom{-}1.128$, $\phantom{-}1.208$] & [$\phantom{-}1.128$, $\phantom{-}1.208$] & never / $>$3 / $>$3 & $>$3 & -- \\
Sebastiani\_2005 & EDN1.3 $\to$ EDNI1.7 & 7 (0) & $-1.212$ & $-1.209$ [$-1.219$, $-1.199$] & [$-1.234$, $-1.179$] & [$-1.234$, $-1.179$] & never / $>$3 / $>$3 & $>$3 & -- \\
Shrier\_2008 & WarmUpExercises $\to$ Injury & 2 (0) & $-1.351$ & $-1.348$ [$-1.358$, $-1.339$] & [$-1.431$, $-1.228$] & [$-1.431$, $-1.228$] & never / $>$3 / $>$3 & never & -- \\
Thoemmes\_2013 & x $\to$ y & 1 (1) & $-0.923$ & $-0.913$ [$-0.928$, $-0.898$] & [$-1.002$, $-0.871$] & [$-1.002$, $-0.871$] & never / $>$3 / $>$3 & never & -- \\
mediator & X $\to$ Y & 3 (3) & $-2.603$ & $-2.600$ [$-2.623$, $-2.576$] & [$-2.623$, $-1.109$] & [$-2.623$, $\phantom{-}0.000$] & 2 / 2 / 2 & 1 & Z$\to$I \,/\, Z$\to$X \\
paths & E $\to$ D & 1 (0) & $\phantom{-}1.194$ & $\phantom{-}1.193$ [$\phantom{-}1.188$, $\phantom{-}1.197$] & [$\phantom{-}1.169$, $\phantom{-}1.216$] & [$\phantom{-}1.169$, $\phantom{-}1.216$] & never / $>$3 / $>$3 & never & -- \\
\addlinespace
\multicolumn{10}{@{}l}{\itshape Networks with coefficients fitted to real data: the frame rows whose analyst graph commits to a set.}\\[1pt]
ecoli70 & asnA $\to$ aceB & 10 (1) & $\phantom{-}0.547$ & $\phantom{-}0.542$ [$\phantom{-}0.511$, $\phantom{-}0.572$] & [$-0.011$, $\phantom{-}0.572$] & [$-0.011$, $\phantom{-}0.572$] & 1 / 1 / 1..$>$3 & 1 & asnA$\to$icdA \\
ecoli70 & b1191 $\to$ yfaD & 10 (1) & $-0.098$ & $-0.097$ [$-0.105$, $-0.089$] & [$-0.114$, $\phantom{-}0.029$] & [$-0.114$, $\phantom{-}0.029$] & 1 / 1 / 1 & 1 & b1191$\to$fixC \\
ecoli70 & cspG $\to$ hupB & 10 (7) & $\phantom{-}0.257$ & $\phantom{-}0.265$ [$\phantom{-}0.246$, $\phantom{-}0.283$] & [$\phantom{-}0.003$, $\phantom{-}0.308$] & [$\phantom{-}0.003$, $\phantom{-}0.308$] & 1 / $>$3 / 1..$>$3 & 1 & -- \\
ecoli70 & pspB $\to$ nmpC & 10 (7) & $-0.108$ & $-0.100$ [$-0.118$, $-0.082$] & [$-0.531$, $\phantom{-}0.026$] & [$-0.531$, $\phantom{-}0.026$] & 1 / 1 / 1..3 & 1 & pspB$\to$pspA \\
ecoli70 & eutG $\to$ lacY & 10 (1) & $\phantom{-}0.309$ & $\phantom{-}0.309$ [$\phantom{-}0.279$, $\phantom{-}0.338$] & [$\phantom{-}0.277$, $\phantom{-}0.377$] & [$\phantom{-}0.277$, $\phantom{-}0.377$] & never / $>$3 / $>$3 & 1 & -- \\
arth150 & 100 $\to$ 245 & 29 (3) & $\phantom{-}1.535$ & $\phantom{-}1.534$ [$\phantom{-}1.527$, $\phantom{-}1.542$] & [$\phantom{-}0.000$, $\phantom{-}1.545$] & [$\phantom{-}0.000$, $\phantom{-}1.545$] & 1 / 1 / 1 & 1 & 100$\to$245 \\
arth150 & 81 $\to$ 464 & 29 (10) & $\phantom{-}0.256$ & $\phantom{-}0.184$ [$\phantom{-}0.118$, $\phantom{-}0.250$] & [$\phantom{-}0.097$, $\phantom{-}0.265$] & [$\phantom{-}0.097$, $\phantom{-}0.265$] & $>$3 / $>$3 / $>$3 & $\geq$2 & -- \\
magic-niab & G1853 $\to$ G43 & 8 (1) & $-0.039$ & $-0.028$ [$-0.038$, $-0.017$] & [$-0.038$, $\phantom{-}0.020$] & [$-0.038$, $\phantom{-}0.020$] & 1 / 1 / 1..$>$3 & 1 & G1853$\to$G311 \\
magic-niab & G311 $\to$ G43 & 8 (1) & $-0.255$ & $-0.251$ [$-0.259$, $-0.244$] & [$-0.260$, $\phantom{-}0.000$] & [$-0.260$, $\phantom{-}0.000$] & 1 / 1 / 1 & 1 & G1853$\to$G311 \\
magic-niab & G418 $\to$ FUS & 8 (2) & $\phantom{-}0.003$ & $-0.003$ [$-0.015$, $\phantom{-}0.010$] & [$-0.022$, $\phantom{-}0.011$] & [$-0.022$, $\phantom{-}0.011$] & 1 / 0 / 0..1 & 1 & -- \\
magic-niab & G1217 $\to$ FT & 8 (1) & $\phantom{-}0.068$ & $\phantom{-}0.015$ [$-0.029$, $\phantom{-}0.060$] & [$-0.031$, $\phantom{-}0.131$] & [$-0.031$, $\phantom{-}0.131$] & never / 0 / 0..$>$3 & 1 & -- \\
magic-niab & G1217 $\to$ YLD & 8 (1) & $\phantom{-}0.009$ & $\phantom{-}0.011$ [$\phantom{-}0.003$, $\phantom{-}0.018$] & [$\phantom{-}0.004$, $\phantom{-}0.021$] & [$\phantom{-}0.004$, $\phantom{-}0.021$] & never / $>$3 / 0..$>$3 & 1 & -- \\
magic-irri & G1378 $\to$ SUB & 10 (1) & $-0.035$ & $-0.031$ [$-0.044$, $-0.017$] & [$-0.048$, $\phantom{-}0.021$] & [$-0.048$, $\phantom{-}0.021$] & 1 / 1 / 1 & 1 & G1378$\to$G1311 \\
magic-irri & G3212 $\to$ FT & 10 (2) & $\phantom{-}0.004$ & $-0.025$ [$-0.070$, $\phantom{-}0.020$] & [$-0.143$, $\phantom{-}0.008$] & [$-0.143$, $\phantom{-}0.008$] & 1 / 0 / 0..1 & 1 & -- \\
magic-irri & G3925 $\to$ GTEMP & 10 (1) & $\phantom{-}0.000$ & $\phantom{-}0.001$ [$-0.034$, $\phantom{-}0.036$] & [$-0.131$, $\phantom{-}0.038$] & [$-0.131$, $\phantom{-}0.038$] & 1 / 0 / 0 & 1 & -- \\
magic-irri & G4145 $\to$ G3222 & 10 (2) & $-0.020$ & $-0.015$ [$-0.025$, $-0.004$] & [$-0.026$, $\phantom{-}0.026$] & [$-0.026$, $\phantom{-}0.026$] & 1 / 0 / 0..1 & 1 & -- \\
magic-irri & G4156 $\to$ CHALK & 10 (2) & $\phantom{-}0.404$ & $\phantom{-}0.501$ [$\phantom{-}0.361$, $\phantom{-}0.641$] & [$-0.154$, $\phantom{-}0.664$] & [$-0.154$, $\phantom{-}0.664$] & 1 / 1 / 1..$>$3 & 1 & G4156$\to$G4145 \\
\bottomrule
\end{tabular}
\end{table}
\end{landscape}


\noindent Excluded, with the reason: M-bias, E on D, because the published diagram is not a
directed acyclic graph; Schipf 2010, TT on T2DM, because the outcome is not a possible descendant of
the treatment, so there is no effect to report; confounding, E on D, because the treatment has no
undirected edge, so no claim of the analyst's bears on it.

\subsection{What the table says once you have read it twice}

The published analyses are robust and the sampled ones are not. Of the nine declared queries that
can be certified, eight are unbroken by any retraction down to depth three: their interval is a
point, their radius is never reached, and no claim the analyst made is load-bearing. The one
exception is mediator, a four-node textbook diagram. Below the rule, on the four networks with
fitted coefficients, sixteen of seventeen informative queries break at the first retraction.

The reason is mechanical and it belongs in the paper. A query enters the sampled frame only if some
retraction changes the validity of the committed set, so a query that no retraction can touch is
rejected before it is measured. The population the method is evaluated on is therefore close to the
complement of the queries these source papers exist to answer, which is an applicability statement
rather than a result.

The radius is almost always one. Across the whole table $r_{\mathrm{val}}$ takes the value one on
every informative row except mediator. In practice the quantity is a pointer to a single claim
rather than a graded scale.

The null radius and the breakdown radius are different numbers and they disagree. On mediator the
set stops being valid after one revision while the interval only reaches zero after two. On several
fitted rows the null radius read from the band at a single seed is zero, meaning the reported
interval already contained zero before any retraction: at that sample size there was no conclusion
to break. The range column exists because that reading moves with the seed, so any sentence about
the null radius has to carry which of the three readings it uses.

One retraction is not a small perturbation. On sixteen of the seventeen informative rows the
interval after retracting one claim is the whole class interval, to four decimals. All the
identifying power the knowledge supplies sits in a single claim, and withdrawing it returns the
analyst to where they would have been with no background knowledge at all. That fact is what the
figure in the next section is built around.
